\documentclass[handout,aspectratio=169,10pt]{beamer}

% =====================================================
% PAQUETES
% =====================================================
\usepackage[utf8]{inputenc}
\usepackage[T1]{fontenc}
\usepackage[spanish,es-noshorthands]{babel}
\usepackage{lmodern}
\usepackage{amsmath,amssymb,mathtools}
\usepackage{array}
\usepackage{booktabs}
\usepackage{multicol}
\usepackage{tikz}
\usepackage{ragged2e}
\usepackage{xcolor}

\usetikzlibrary{calc,positioning,shapes.geometric}

% =====================================================
% TEMA
% =====================================================
\usetheme{Madrid}
\usecolortheme{default}
\setbeamertemplate{navigation symbols}{}

% =====================================================
% COLORES UST MODERNOS
% =====================================================
\definecolor{USTBlue}{RGB}{0,70,127}
\definecolor{USTBlueDark}{RGB}{0,45,85}
\definecolor{USTTeal}{RGB}{0,153,117}
\definecolor{USTGray}{RGB}{95,95,95}
\definecolor{USTLight}{RGB}{245,247,250}
\definecolor{USTSoft}{RGB}{232,239,247}
\definecolor{USTText}{RGB}{35,40,45}
\definecolor{USTAlert}{RGB}{180,50,47}
\definecolor{USTGold}{RGB}{201,162,39}

% =====================================================
% CONFIGURACIÓN VISUAL
% =====================================================
\setbeamercolor{normal text}{fg=USTText,bg=white}
\setbeamercolor{structure}{fg=USTBlue}
\setbeamercolor{frametitle}{bg=USTBlue,fg=white}
\setbeamercolor{title}{fg=white}
\setbeamercolor{subtitle}{fg=white}

\setbeamercolor{block title}{bg=USTBlue,fg=white}
\setbeamercolor{block body}{bg=USTLight,fg=black}

\setbeamercolor{block title alerted}{bg=USTAlert,fg=white}
\setbeamercolor{block body alerted}{bg=red!5,fg=black}

\setbeamercolor{block title example}{bg=USTTeal,fg=white}
\setbeamercolor{block body example}{bg=green!5,fg=black}

\setbeamertemplate{blocks}[rounded][shadow=false]

\setbeamerfont{title}{series=\bfseries,size=\LARGE}
\setbeamerfont{subtitle}{size=\large}
\setbeamerfont{frametitle}{series=\bfseries,size=\large}
\setbeamerfont{block title}{series=\bfseries}

% =====================================================
% FOOTLINE PERSONALIZADA
% =====================================================
\setbeamertemplate{footline}{
	\leavevmode
	\hbox{
		\begin{beamercolorbox}[wd=.72\paperwidth,ht=2.8ex,dp=1ex,leftskip=1em]{author in head/foot}
			\color{USTBlueDark}\small\textbf{Pensamiento Matemático} \hspace{0.5em} | \hspace{0.5em}
			Unidad II – Operatoria Matemática Básica Aplicada a las Ciencias Sociales
		\end{beamercolorbox}
		\begin{beamercolorbox}[wd=.28\paperwidth,ht=2.8ex,dp=1ex,rightskip=1em plus1fil]{date in head/foot}
			\color{USTBlueDark}\small\insertframenumber/\inserttotalframenumber
		\end{beamercolorbox}
	}
	\vskip0pt
}

% =====================================================
% ENCABEZADO CON BARRA
% =====================================================
\setbeamertemplate{headline}{
	\begin{tikzpicture}[remember picture,overlay]
		\fill[USTBlue] (current page.north west) rectangle ([yshift=-0.18cm]current page.north east);
		\fill[USTTeal] (current page.north west) rectangle ([xshift=3.2cm,yshift=-0.18cm]current page.north west);
	\end{tikzpicture}
	\vspace{0.15cm}
}

% =====================================================
% COMANDOS PERSONALES
% =====================================================
\newcommand{\IdeaClave}[1]{
	\begin{exampleblock}{Idea clave}
		#1
	\end{exampleblock}
}

\newcommand{\AlertaPedagogica}[1]{
	\begin{alertblock}{Atención}
		#1
	\end{alertblock}
}

\newcommand{\ObjetivoClase}[1]{
	\begin{block}{Objetivo de la clase}
		#1
	\end{block}
}

% =====================================================
% DATOS
% =====================================================
\title{Pensamiento Matemático}
\subtitle{Unidad II – Operatoria Matemática Básica Aplicada a las Ciencias Sociales}
\author{Profesor Luis González Alcaino}
\institute{Universidad Santo Tomás \\ Departamento Ciencias Básicas - Talca}
\date{2026}

% =====================================================
% DOCUMENTO
% =====================================================
\begin{document}
	
	% =====================================================
	% PORTADA
	% =====================================================
	\begin{frame}[plain]
		\begin{tikzpicture}[remember picture,overlay]
			\fill[USTBlue] (current page.south west) rectangle (current page.north east);
			\fill[USTBlueDark] (current page.south west) rectangle ([xshift=3.2cm]current page.north west);
			\fill[USTTeal] ([yshift=-1.2cm]current page.north west) rectangle ([yshift=-1.5cm]current page.north east);
			\fill[white,opacity=0.08] ([xshift=9cm,yshift=-1cm]current page.north west) circle (1.7cm);
			\fill[white,opacity=0.05] ([xshift=11.5cm,yshift=-4cm]current page.north west) circle (2.3cm);
		\end{tikzpicture}
		
		\vspace{1.2cm}
		
		{\color{white}\bfseries\LARGE \inserttitle\par}
		\vspace{0.3cm}
		{\color{white}\Large \insertsubtitle\par}
		\vspace{0.9cm}
		
		\begin{beamercolorbox}[wd=\textwidth,rounded=true,sep=0.8em]{block body}
			\centering
			\textbf{\insertauthor}\\
			\insertinstitute\\
			\insertdate
		\end{beamercolorbox}
		
		\vfill
		
		{\color{white!85}\small Clase: Conjuntos Numéricos y Operatoria}
	\end{frame}
	
	% =====================================================
	% RUTA DE LA CLASE
	% =====================================================
	\begin{frame}{Ruta de la clase}
		\tableofcontents
	\end{frame}
	
	% =====================================================
	\section{Resultados de Aprendizaje}
	% =====================================================
	
	\begin{frame}{Resultados de Aprendizaje}
		\begin{itemize}[<+->]
			\item Aplica operatoria con números enteros y racionales en la resolución de ejercicios y situaciones contextualizadas.
			\item Integra patrones y relaciones matemáticas al analizar propiedades, estructuras y resultados en distintas expresiones numéricas.
			\item Detecta patrones, relaciones e inconsistencias en procedimientos y resultados, justificando sus conclusiones mediante razonamiento matemático.
		\end{itemize}
	\end{frame}
	
	\begin{frame}{¿Qué aprenderemos hoy?}
		\ObjetivoClase{
			Resolver operaciones con números enteros y racionales, reconociendo propiedades, patrones, relaciones e inconsistencias en distintos procedimientos matemáticos vinculados al análisis cuantitativo en contextos de las ciencias sociales.
		}
		
		\IdeaClave{
			No solo importa obtener un resultado correcto: también es fundamental reconocer cómo se construye, qué patrón sigue y cuándo un procedimiento contiene un error.
		}
	\end{frame}
	
	% =====================================================
	\section{Conjuntos Numéricos}
	% =====================================================
	
	\begin{frame}{Conjuntos Numéricos}
		\begin{block}{Introducción}
			Se estudiarán los conjuntos numéricos que conforman a los números reales, vinculándolos con ejemplos frecuentes en medición, conteo y análisis de datos en contextos psicológicos y sociales.
		\end{block}
		
		\begin{itemize}[<+->]
			\item Números naturales \(\mathbb{N}\)
			\item Números enteros \(\mathbb{Z}\)
			\item Números racionales \(\mathbb{Q}\)
			\item Números irracionales \(\mathbb{I}\)
			\item Números reales \(\mathbb{R}\)
		\end{itemize}
	\end{frame}
	
	\begin{frame}{Los números naturales \(\mathbb{N}\)}
		\begin{block}{Definición}
			También conocidos como números para contar o enteros positivos.
			\[
			\mathbb{N}=\{1,2,3,\dots\}
			\]
		\end{block}
		
		\begin{exampleblock}{Contexto}
			En Psicología o Ciencias Sociales pueden utilizarse para contar:
			\begin{itemize}
				\item participantes de un estudio,
				\item sesiones de intervención,
				\item respuestas correctas en un test,
				\item frecuencia de conductas observadas.
			\end{itemize}
		\end{exampleblock}
	\end{frame}
	
	\begin{frame}{Los números enteros \(\mathbb{Z}\)}
		\begin{block}{Definición}
			Permiten representar cantidades positivas, negativas y el cero.
			\[
			\mathbb{Z}=\{\dots,-3,-2,-1,0,1,2,3,\dots\}
			\]
		\end{block}
		
		\begin{exampleblock}{Contexto}
			Son útiles para representar:
			\begin{itemize}
				\item aumentos o disminuciones de puntajes,
				\item variaciones en escalas de bienestar,
				\item cambios en niveles de ansiedad o motivación.
			\end{itemize}
		\end{exampleblock}
	\end{frame}
	
	\begin{frame}{Los números racionales \(\mathbb{Q}\)}
		\begin{block}{Definición}
			Los números racionales son aquellos que pueden escribirse como cociente de dos enteros:
			\[
			\mathbb{Q}=\left\{\frac{a}{b}: a,b\in\mathbb{Z},\ b\neq 0\right\}
			\]
		\end{block}
		
		\begin{exampleblock}{Ejemplos}
			\[
			\frac14=0.25,\qquad \frac53=1.666\dots,\qquad \frac{19}{12}=1.58333\dots
			\]
		\end{exampleblock}
		
		\begin{block}{Aplicación}
			Aparecen en porcentajes, proporciones, promedios y razones entre variables observadas.
		\end{block}
	\end{frame}
	
	\begin{frame}{Los números irracionales \(\mathbb{I}\)}
		\begin{block}{Definición}
			Son números que no pueden expresarse como cociente de enteros y tienen desarrollo decimal infinito no periódico.
			\[
			\mathbb{I}=\{x:\ x\notin\mathbb{Q}\}
			\]
		\end{block}
		
		\begin{exampleblock}{Ejemplos}
			\[
			\sqrt{2}=1.414213\dots,\qquad \pi=3.141592\dots,\qquad e=2.718281\dots
			\]
		\end{exampleblock}
	\end{frame}
	
	\begin{frame}{Los números reales \(\mathbb{R}\)}
		\begin{block}{Definición}
			El conjunto de los números reales está formado por la unión de los racionales y los irracionales:
			\[
			\mathbb{R}=\mathbb{Q}\cup\mathbb{I}
			\]
		\end{block}
		
		\IdeaClave{
			Los números reales permiten representar medidas, escalas, puntajes, diferencias y comparaciones cuantitativas en una amplia variedad de contextos.
		}
	\end{frame}
	
	\begin{frame}{Recta numérica y comparación}
		\begin{center}
			\begin{tikzpicture}[scale=0.9]
				\draw[->,thick] (-4,0)--(4.5,0);
				\foreach \x in {-3,-2,-1,0,1,2,3,4}
				{
					\draw (\x,0.12)--(\x,-0.12);
					\node[below] at (\x,-0.12) {\x};
				}
			\end{tikzpicture}
		\end{center}
		
		\begin{block}{Comparación}
			\begin{itemize}[<+->]
				\item \(7<12\)
				\item \(-3>-10\)
				\item \(10\leq 15\)
				\item \(8\geq 8\)
			\end{itemize}
		\end{block}
		
		\begin{exampleblock}{Interpretación}
			Comparar números permite analizar diferencias entre puntajes, resultados observados o niveles en una escala.
		\end{exampleblock}
	\end{frame}
	
	\begin{frame}{Ejercicio: clasificación en conjuntos numéricos}
		\begin{block}{Instrucción}
			Marque con \(\checkmark\) el o los conjuntos numéricos a los cuales pertenece cada número.
		\end{block}
		
		\[
		\begin{array}{c|ccccc}
			\text{Número} & \mathbb{N} & \mathbb{Z} & \mathbb{Q} & \mathbb{I} & \mathbb{R} \\ \hline
			-5 & & & & & \\
			\sqrt{3} & & & & & \\
			\frac{5}{8} & & & & & \\
			\pi & & & & & \\
			2 & & & & &
		\end{array}
		\]
	\end{frame}
	
	% =====================================================
	\section{Operatoria con Números Enteros}
	% =====================================================
	
	\begin{frame}{Adición de números enteros}
		\begin{block}{Reglas básicas}
			\begin{enumerate}[<+->]
				\item Si tienen el mismo signo, se suman los valores absolutos y se conserva el signo.
				\item Si tienen distinto signo, se restan los valores absolutos y se conserva el signo del número con mayor valor absoluto.
			\end{enumerate}
		\end{block}
		
		\begin{exampleblock}{Ejemplos}
			\[
			2+5=7,\qquad -4+(-7)=-11,\qquad 8+(-3)=5,\qquad -10+9=-1
			\]
		\end{exampleblock}
	\end{frame}
	
	\begin{frame}{Observaciones sobre la adición}
		\begin{itemize}[<+->]
			\item \(a+(-b)=a-b\)
			\item \(a-(-b)=a+b\)
			\item \(a+b=b+a\) \hfill (conmutativa)
			\item \(a+(b+c)=(a+b)+c\) \hfill (asociativa)
			\item \(a+0=a\)
			\item \(a+(-a)=0\)
			\item \(-(-a)=a\)
			\item \(-(a+b)=(-a)+(-b)\)
		\end{itemize}
	\end{frame}
	
	\begin{frame}{Ejercicios con enteros}
		\begin{block}{Determine el valor de las siguientes expresiones}
			\begin{multicols}{2}
				\begin{enumerate}[a)]
					\item \((-3)+(-8)\)
					\item \((-5+3)+2\)
					\item \(-10+5\)
					\item \(0+(-5+2)\)
					\item \((-32)+48+(-10)\)
					\item \(-(-14)+(-18)\)
					\item \(-(3+(-7))-7-3\)
					\item \(-(-8+5)+(-8)\)
				\end{enumerate}
			\end{multicols}
		\end{block}
	\end{frame}
	
	\begin{frame}{Multiplicación de números enteros}
		\begin{block}{Ley de signos}
			\[
			(+)\cdot(+)=+,\qquad (-)\cdot(-)=+,\qquad (+)\cdot(-)=-,\qquad (-)\cdot(+)=-
			\]
		\end{block}
		
		\begin{exampleblock}{Ejemplos}
			\[
			(-2)\cdot 7=-14,\qquad (-3)\cdot(-7)=21,\qquad 4\cdot(-5)=-20
			\]
		\end{exampleblock}
	\end{frame}
	
	\begin{frame}{Propiedades de la multiplicación}
		\begin{itemize}[<+->]
			\item \(a\cdot b=b\cdot a\) \hfill (conmutativa)
			\item \(a\cdot(b\cdot c)=(a\cdot b)\cdot c\) \hfill (asociativa)
			\item \(a\cdot 1=a\)
			\item \(a\cdot(b+c)=a\cdot b+a\cdot c\) \hfill (distributiva)
			\item \(a\cdot 0=0\)
			\item Si \(a\cdot b=0\), entonces \(a=0\) o \(b=0\)
		\end{itemize}
	\end{frame}
	
	\begin{frame}{Ejercicios de multiplicación con enteros}
		\begin{multicols}{2}
			\begin{enumerate}[a)]
				\item \((-5)\cdot 0\)
				\item \(((-10)\cdot 4)\cdot 5\)
				\item \(((-20)+40)\cdot(-2)\)
				\item \(4\cdot(-5)\)
				\item \(20\cdot((-8)-6)\)
				\item \((-3)\cdot(-7)\)
				\item \(3+2\cdot 5\)
				\item \((-2)\cdot(-(-14)+(-5))\)
				\item \((-7+3)\cdot(-3)\)
				\item \(5\cdot4-5+6\cdot3-2+7\cdot3\)
			\end{enumerate}
		\end{multicols}
	\end{frame}
	
	\begin{frame}{Jerarquía de operaciones}
		\begin{block}{Orden de resolución}
			\begin{enumerate}[<+->]
				\item Resolver paréntesis.
				\item Efectuar multiplicaciones y divisiones.
				\item Realizar sumas y restas.
				\item Si hay operaciones de igual prioridad, proceder de izquierda a derecha.
			\end{enumerate}
		\end{block}
		
		\begin{exampleblock}{Ejemplos}
			\[
			2+[3-(5-2\cdot3)\cdot4]
			\]
			\[
			(5\cdot3-2)-[4\cdot(5\cdot3)-1]
			\]
		\end{exampleblock}
	\end{frame}
	
	% =====================================================
	\section{Operatoria con Números Racionales}
	% =====================================================
	
	\begin{frame}{Adición de números racionales}
		\begin{block}{Reglas}
			\begin{enumerate}[<+->]
				\item Si tienen igual denominador:
				\[
				\frac{a}{b}+\frac{c}{b}=\frac{a+c}{b},\qquad b\neq0
				\]
				\item Si tienen distinto denominador:
				\[
				\frac{a}{b}+\frac{c}{d}=\frac{ad+bc}{bd},\qquad b\neq0,\ d\neq0
				\]
			\end{enumerate}
		\end{block}
	\end{frame}
	
	\begin{frame}{Ejemplos con racionales}
		\[
		\frac12-\frac52+\frac32=\frac{1-5+3}{2}=-\frac12
		\]
		
		\vspace{0.3cm}
		
		\[
		\frac13-\frac25=\frac{5-6}{15}=-\frac1{15}
		\]
		
		\vspace{0.3cm}
		
		\[
		\frac45+\frac83-\frac7{18}+\frac92=\frac{341}{45}
		\]
		
		\vspace{0.3cm}
		
		\[
		\frac14+\frac75-\frac38=\frac{51}{40}
		\]
	\end{frame}
	
	\begin{frame}{Observaciones sobre racionales}
		\begin{itemize}[<+->]
			\item \(\displaystyle \frac{-a}{b}=\frac{a}{-b}=-\frac{a}{b}\)
			\item \(\displaystyle \frac{-a}{-b}=\frac{a}{b}\)
			\item \(\displaystyle -\left(\frac{a}{b}\right)=-\frac{a}{b}\)
			\item \(\displaystyle \frac{a}{1}=a\)
			\item \(\displaystyle \frac{a}{a}=1,\ a\neq0\)
			\item \(\displaystyle \frac{0}{a}=0,\ a\neq0\)
			\item \(\displaystyle \frac{a}{0}\) no existe
		\end{itemize}
	\end{frame}
	
	\begin{frame}{Multiplicación de racionales}
		\begin{block}{Regla}
			\[
			\frac{a}{b}\cdot\frac{c}{d}=\frac{ac}{bd},\qquad b\neq0,\ d\neq0
			\]
		\end{block}
		
		\begin{exampleblock}{Ejemplos}
			\[
			\frac15\cdot\left(-\frac38\right)=-\frac3{40}
			\]
			\[
			\frac3{-8}\cdot\frac57=-\frac{15}{56}
			\]
			\[
			6\cdot\frac{5}{11}=\frac{30}{11}
			\]
		\end{exampleblock}
	\end{frame}
	
	\begin{frame}{Propiedades y equivalencia de fracciones}
		\begin{itemize}[<+->]
			\item \(\displaystyle \frac{a}{b}=\frac{c}{d}\iff a\cdot d=b\cdot c\)
			\item Simplificar una fracción consiste en cancelar un factor común.
			\item Amplificar una fracción consiste en multiplicar numerador y denominador por un mismo número.
		\end{itemize}
		
		\begin{exampleblock}{Ejemplos}
			\[
			\frac{8}{12}=\frac{2\cdot4}{3\cdot4}=\frac23
			\qquad\qquad
			-\frac38=\frac{-3\cdot4}{8\cdot4}=-\frac{12}{32}
			\]
		\end{exampleblock}
	\end{frame}
	
	\begin{frame}{Ejercicios con multiplicación de racionales}
		\begin{multicols}{2}
			\begin{enumerate}[a)]
				\item \(\frac39\cdot\left(-\frac43\right)\)
				\item \(\left((-2)-\frac13\right)\cdot\frac23\)
				\item \(\left(-\frac63\right)\cdot\left(-\frac47\right)\cdot\frac19\)
				\item \(\frac27\cdot\frac12\cdot(-3)\)
				\item \((-2)\cdot\frac6{-3}\cdot\frac{(-12)(-9)}4\)
				\item \(\left(-\frac78\right)\cdot\left(\frac{-4}{-3}\right)\cdot\frac{2}{14}\cdot(-4)\)
			\end{enumerate}
		\end{multicols}
	\end{frame}
	
	\begin{frame}{División de racionales}
		\begin{block}{Regla}
			Dividir por una fracción equivale a multiplicar por su inversa:
			\[
			\frac{a}{b}\div\frac{c}{d}=\frac{a}{b}\cdot\frac{d}{c}
			\]
		\end{block}
		
		\begin{exampleblock}{Ejemplos}
			\[
			\frac35\div\frac29=\frac35\cdot\frac92=\frac{27}{10}
			\]
			\[
			\frac12\div\left(-\frac23\right)=\frac12\cdot\frac3{-2}=-\frac34
			\]
			\[
			7\div\frac{14}{5}=7\cdot\frac5{14}=\frac52
			\]
		\end{exampleblock}
	\end{frame}
	
	\begin{frame}{Ejercicios de división de racionales}
		\begin{multicols}{2}
			\begin{enumerate}[a)]
				\item \(-\frac85\div\frac47\)
				\item \(\left(-\frac13\right)\div\left(\frac2{-3}\right)\)
				\item \(\left(-\frac34\right)\div 6\)
				\item \(\frac79\div\left(-\frac4{12}\right)\)
				\item \(\left(\frac{(-2)\cdot 6}{-3}\right)\div\left(\frac{(-15)(-2)}4\right)\)
				\item \(\left(\left(-\frac58\right)\cdot\left(\frac{-4}{-5}\right)\right)\div\left(\frac{12}{6}\right)\)
			\end{enumerate}
		\end{multicols}
	\end{frame}
	
	% =====================================================
	\section{Patrones, Relaciones e Inconsistencias}
	% =====================================================
	
	\begin{frame}{Actividad: Reconociendo Patrones}
		\begin{block}{Instrucción}
			Analice las siguientes expresiones y determine qué patrón, propiedad o relación matemática se está utilizando.
		\end{block}
		
		\begin{enumerate}[<+->]
			\item \((-3)+(-8)=-11\)
			\item \(8+(-3)=5\)
			\item \((-25)+10=10+(-25)\)
			\item \(((-12)+15)+(-3)=(-12)+(15-3)\)
		\end{enumerate}
		
		\begin{alertblock}{Pregunta de análisis}
			¿Qué tienen en común estas expresiones? ¿Qué propiedades permiten justificar cada paso?
		\end{alertblock}
	\end{frame}
	
	\begin{frame}{Actividad: Detectar Errores}
		\begin{block}{Instrucción}
			Analice los siguientes procedimientos. Determine si son correctos. Si existe un error, identifíquelo y corríjalo.
		\end{block}
		
		\begin{enumerate}[<+->]
			\item \(-4+(-7)=3\)
			\item \(6+(-8)=6+8=14\)
			\item \((-2)\cdot(-3)=-6\)
			\item \(\frac12+\frac13=\frac25\)
		\end{enumerate}
		
		\AlertaPedagogica{
			Un resultado incorrecto puede originarse por una mala aplicación de signos, por ignorar la jerarquía de operaciones o por operar fracciones de forma incorrecta.
		}
	\end{frame}
	
	\begin{frame}{Relaciones matemáticas que debemos reconocer}
		\begin{columns}[T]
			\begin{column}{0.48\textwidth}
				\begin{exampleblock}{Relaciones}
					\begin{itemize}[<+->]
						\item igualdad,
						\item desigualdad,
						\item equivalencia,
						\item opuesto aditivo,
						\item inverso multiplicativo.
					\end{itemize}
				\end{exampleblock}
			\end{column}
			
			\begin{column}{0.48\textwidth}
				\begin{block}{Patrones frecuentes}
					\begin{itemize}[<+->]
						\item repetición de reglas de signo,
						\item conservación de propiedades,
						\item uso sistemático del denominador común,
						\item regularidad en transformaciones equivalentes.
					\end{itemize}
				\end{block}
			\end{column}
		\end{columns}
	\end{frame}
	
	% =====================================================
	\section{Aplicaciones}
	% =====================================================
	
	\begin{frame}{Aplicación 1: Promedio de estrés}
		\begin{block}{Situación}
			Un psicólogo realiza un estudio sobre el estrés en 10 pacientes. Los niveles reportados son:
			\[
			4,\ 5,\ 6,\ 3,\ 8,\ 7,\ 5,\ 6,\ 9,\ 10
			\]
		\end{block}
		
		\begin{enumerate}[<+->]
			\item Calcule la suma total de los puntajes.
			\item Determine el promedio del índice de estrés.
			\item Interprete el resultado.
		\end{enumerate}
	\end{frame}
	
	\begin{frame}{Aplicación 2: Reducción promedio de ansiedad}
		\begin{block}{Situación}
			Se registran niveles de ansiedad antes y después de un tratamiento en 5 pacientes.
			
			\[
			\text{Antes: } 8,\ 7,\ 9,\ 6,\ 10
			\]
			\[
			\text{Después: } 5,\ 4,\ 6,\ 3,\ 7
			\]
		\end{block}
		
		\begin{enumerate}[<+->]
			\item Determine la reducción en cada paciente.
			\item Calcule la reducción promedio.
			\item Analice si existe un patrón común en las variaciones observadas.
		\end{enumerate}
	\end{frame}
	
	\begin{frame}{Aplicación 3: Puntaje Z}
		\begin{block}{Situación}
			Un paciente obtiene:
			\[
			X_A=85,\quad \mu_A=70,\quad \sigma_A=10
			\]
			\[
			X_B=90,\quad \mu_B=75,\quad \sigma_B=15
			\]
			Se utiliza la fórmula:
			\[
			Z=\frac{X-\mu}{\sigma}
			\]
		\end{block}
		
		\begin{enumerate}[<+->]
			\item Calcule el puntaje Z en la prueba A.
			\item Calcule el puntaje Z en la prueba B.
			\item Compare ambos resultados.
		\end{enumerate}
	\end{frame}
	
	\begin{frame}{Aplicación integradora en contexto psicológico}
		\begin{block}{Situación}
			Durante una semana, un paciente presenta las siguientes variaciones en su nivel de bienestar:
			\[
			+5,\ -3,\ +2,\ -6,\ +4
			\]
		\end{block}
		
		\begin{enumerate}[<+->]
			\item Determine el cambio total en el bienestar.
			\item Identifique patrones en las variaciones.
			\item Detecte si existe alguna inconsistencia al interpretar los resultados.
			\item Explique qué conclusión podría extraerse desde el contexto psicológico.
		\end{enumerate}
	\end{frame}
	
	% =====================================================
	\section{Cierre}
	% =====================================================
	
	\begin{frame}{Síntesis de la clase}
		\begin{block}{Hoy trabajamos}
			\begin{itemize}[<+->]
				\item Clasificación de los conjuntos numéricos.
				\item Operatoria con números enteros y racionales.
				\item Propiedades y jerarquía de operaciones.
				\item Reconocimiento de patrones y relaciones.
				\item Detección de errores e inconsistencias.
				\item Aplicaciones en contextos psicológicos y sociales.
			\end{itemize}
		\end{block}
	\end{frame}
	
	\begin{frame}{Verificación de los aprendizajes}
		\begin{enumerate}[<+->]
			\item ¿Puedo resolver correctamente operaciones con enteros y racionales?
			\item ¿Reconozco patrones y propiedades en los procedimientos?
			\item ¿Puedo detectar cuándo un resultado o una operación es inconsistente?
			\item ¿Soy capaz de interpretar matemáticamente una situación contextualizada?
		\end{enumerate}
	\end{frame}
	
	\begin{frame}[plain]
		\begin{tikzpicture}[remember picture,overlay]
			\fill[USTBlue] (current page.south west) rectangle (current page.north east);
			\fill[USTTeal] ([yshift=-1.1cm]current page.north west) rectangle ([yshift=-1.45cm]current page.north east);
		\end{tikzpicture}
		
		\vfill
		\begin{center}
			{\color{white}\bfseries\Huge Gracias}\\[0.4cm]
			{\color{white!90}\Large Clase 7 – Conjuntos Numéricos y Operatoria}
		\end{center}
		\vfill
	\end{frame}
	
\end{document}